\section{Interceptação e Emulação}

Do inglês \textit{trap-and-emulate}, interceptação e emulação trata-se de um método utilizado para viabilizar a virtualização em sistemas nos quais o suporte de
hardware não é completamente implementado. Nesse modelo, quando a aplicação executada em ambiente virtualizado tenta realizar uma operação que depende de recursos privilegiados,
essa operação é interceptada e tratada como uma exceção. O controle é então transferido para um nível mais privilégiado, geralmente o hipervisor, que fica responsável por
emular em software o comportamento esperado da instrução ou operação originalmente solicitada.

\begin{figure}[H]
\centering
\includegraphics[width=0.8\textwidth]{images/TrapAndEmulate.pdf}
\caption{Fluxograma de Interceptação de Instruções Privilegiadas}\label{fig:trap-and-emulate}
\end{figure}


